\documentclass[10pt]{article}
\usepackage{tmlr}
\usepackage{amsmath,amssymb,booktabs,longtable,array,calc}
\usepackage{graphicx,xcolor,tikz}
\usetikzlibrary{arrows.meta,positioning}
\usepackage{hyperref,url}
\newcommand{\tightlist}{\setlength{\itemsep}{0pt}\setlength{\parskip}{0pt}}
\hypersetup{hidelinks}

\title{One Additional Observation, More Conditional Certificates:\\A Paired Study in a Finite System Class}
\author{Anonymous authors}
\begin{document}\raggedbottom\maketitle
\begingroup\renewcommand{\thefootnote}{}\footnotetext{AI assistance: An LLM-based assistant contributed to experimental design, implementation, analysis, literature review, and manuscript drafting and revision. AI tools are not authors.}\endgroup
\begin{abstract}
Does a development-selected query policy certify more future targets than a
fixed query at the same environmental observation budget? We compare two
deterministic policies on 5,000 pre-specified paired draws from a finite frame
of system classes and histories. A certificate requires unanimous target
predictions in a nonempty compatible candidate set. The selected policy
certifies 372 targets versus 283 for the fixed query: 90 selected-only and
one fixed-only outcome. The net difference is 1.78 percentage points, with
the pre-specified at-least-95\% interval [1.23, 2.29]. All 655 issued certificates
match the privileged audit. Equal observation budgets require unequal local
computation; 92.56\% of selected-policy draws remain uncertified. Coverage of
the true system by the candidate class is a premise. This is a conditional
acquisition result, not a new active-learning principle, live LLM result,
universal guarantee or proof of a gain of at least two percentage points.
\end{abstract}
\section{Question, contract and findings}
\label{observation-selection}

This study asks whether
\emph{acquiring} one further legal observation helps resolve an
uncertainty: a future target outcome left ambiguous by the available
history. It is a separate experiment with deterministic actors, a
covering class of deterministic systems with at most three states and
binary inputs and outputs, and two permitted queries. A conditional
target certificate is issued only when every candidate consistent with
the history and the observed query response predicts the same target.
Class coverage is an explicit premise of this certificate. From the
eight-input history's cutoff, $q_0$ applies input 0 and $q_1$ applies
input 1; each returns only that transition's output. The target is the
pair of second outputs on input sequences $(0,0)$ and $(1,0)$, each
starting at the original cutoff. Appendix~\ref{app:observation-selection}
defines the state and output conventions.

The fixed policy always chooses the first query. The development-selected
policy maximizes the number of nonempty response branches having a
single target outcome, breaks ties by minimizing the largest feasible
branch's number of distinct targets, and finally chooses the first
query. Branch and candidate counts are not calibrated probabilities.
Both arms receive exactly one observation on every draw, including
initially unambiguous and duplicate records. Separate processes and
state copies prevent cross-arm reuse. Neither actor receives the true
system's identity, private state or target answer. All predictions were
committed before a separate privileged target audit.

The protocol fixed the policy, tie rules, $N=5{,}000$, confidence
procedure and validity gate before the main draw. These are local, hash-bound
precommitments, not an independently timestamped public preregistration.
The hypothesis class covers at most three states; the sampled worlds instead
have exactly three minimal reachable states, after development exclusions.
Sampling was uniform
with replacement over 8,372 eligible behavior classes crossed with
256 input histories: 2,143,232 class--history pairs. Classes used in the
specified development samples were excluded; visible-history overlap
was permitted. The realized sample contains 4,993 distinct pairs and
3,791 distinct classes. Draws, rather than unique worlds or live model
attempts, are the inferential units. Appendix~\ref{app:observation-selection}
records the development history, sampling and statistical procedure.

\begin{table}[t]
\centering
\small
\begin{tabular}{lrrr}
\toprule
 & Fixed: certificate & Fixed: no certificate & Total\\
\midrule
Selected: certificate & 282 & 90 & 372\\
Selected: no certificate & 1 & 4,627 & 4,628\\
\midrule
Total & 283 & 4,717 & 5,000\\
\bottomrule
\end{tabular}
\caption{Separate deterministic acquisition study: certificate issuance
on 5,000 paired draws, not P/L/U feasibility labels. All 655 issued
certificates matched their targets; the 282 jointly certified draws
contribute two certificates each. That cell includes 97 draws whose
targets were already unique before the observation. The remaining
4,903 draws yield 186 new certificates in the fixed arm and 275 in
the selected arm. The net gain is $90-1=89$ on the primary $N=5{,}000$.
No certificate means unresolved under the corresponding policy, not an
impossibility proof for every one-observation method.}
\label{tab:observation-selection}
\end{table}


All 5,000 pairs completed, with no technical missingness or unknown work
counters. The selected policy issued 372 certificates (7.44\%), versus
283 (5.66\%) for the fixed policy. There were 90 selected-only and one
fixed-only certificates: 89 net additional certificates, or
\textbf{1.78 percentage points}. The pre-specified, at-least-95\%
confidence interval is \textbf{[1.23, 2.29] percentage points}.
It excludes zero and supports the narrow positive-difference claim
in this frame. It includes the two-percentage-point research-interest
threshold, so a gain of at least that size remains unresolved. The
planning power calculation is not a guarantee or a post hoc estimate
of achieved power. Of the 5,000 draws, 97 already had a unique target
before the extra observation and remain in both arms' totals. Among the
4,903 initially ambiguous draws, the observation newly certified 186
targets in the fixed arm and 275 in the selected arm; their net difference
is again 89. The primary denominator remains all 5,000 draws.

All 655 issued certificates matched the target audit; any contradiction
would have invalidated the certificate result regardless of its
interval. Nevertheless, 4,628 selected-policy draws (92.56\%) remain
uncertified. The one fixed-only certificate also rules out per-record
dominance. Query choices differ on 686 draws. Of 458 draws with a
pre-observation possible gain, only 90 realize it; the possible-loss
draw realizes one loss. These counts separate a favorable candidate
branch from the actual response. They are descriptive diagnostics,
not further primary tests or calibrated response probabilities.

Equal observation budgets did not imply equal computation. Each arm
made 5,000 environment queries and 233,280,000 candidate--history
checks. Relative to the fixed arm, the selected policy made 5,283,246
additional query--candidate evaluations and 20,000 selection-branch
checks: approximately 59,362 and 225, respectively, per \emph{net}
additional certificate. These work categories are reported separately,
not added as if they were disjoint CPU instructions. Within-arm reuse
is recorded separately and is not counted as additional work.
Actor CPU time was 136.38 versus 137.96 seconds; the fixed arm always
ran first, so timing is descriptive for this implementation and host.
No economic-utility or deployment conclusion is drawn.

This realizes the acquisition benefit bounded by the earlier branch scan.
The earlier failed prediction criteria (Appendix A.1) and the separate
decision-record crossing results remain unchanged.
No general crossing benefit, improved LLM experimentation, universal coverage
or cross-contract transfer follows.

\paragraph{Post hoc query comparators on the same inputs.}
A later reader analysis reconstructed both allowed responses on the original
5,000 draws. Always querying $q_1$ yields 219 certificates, versus 283 for
the pre-specified $q_0$ and 372 for the selected policy. A policy choosing
each query with probability $1/2$ independently of the record has an expected
$(283+219)/2=251$ certificates on these inputs. This is an exact expectation,
not an observed randomized run. Reconstructed $q_0$/$q_1$ certificates have
zero target disagreements. These comparisons are retrospective and use no
new sample or live observation; they neither replace the primary 89-net
gain and its interval nor establish optimality among other selectors.
The reader replay computes them from the reconstructed histories, candidate
supports and original world tables, separately from historical work counters.

\section{Position and scope}
Informative query selection is established \citep{cohn1994active}, including
identifiability-directed acquisition \citep{shetty2026idea}. In active automata
learning, Angluin's membership and equivalence-query setting aims to identify
a regular language \citep{angluin1987learning}; here the known covering class
is a finite transducer class, there is no equivalence oracle, and one of two
one-step queries serves a particular future target rather than full system
identification. Query by committee uses disagreement to select informative
queries \citep{seung1992query}. Our frozen branch-purity selector neither
samples a committee nor interprets candidate multiplicities as probabilities;
we make no learning-rate comparison to that method. Our contribution
is this frozen policy comparison, with realized branches, incorrect-certificate
checks and compute accounting. Neither query selection nor unanimity-based
inference in a covering finite class is claimed as novel.

The positive average difference is limited to the pre-specified finite frame;
new classes, query contracts and live LLM selection remain unassessed. The
primary comparison establishes superiority to the frozen fixed $q_0$
baseline. The $q_1$ and fair-random expectation above are separate post hoc
descriptions, not pre-specified confirmatory tests of a wider optimality claim.
The reader replay reconstructs the original sample, histories, query responses,
candidate computations and target checks from disclosed synthetic inputs,
then recovers the counts and primary interval. It collects no new sample
and does not independently attest the historical run. This result and the crossing
case study retain separate denominators and claims.

\bibliography{references}\bibliographystyle{tmlr}
\appendix\section{Separate finite-class observation-selection study}
\label{app:observation-selection}

\subsection{Development history and a changed estimand}

Earlier prediction criteria remain failed. In one 1,600-draw study,
candidate consensus issued 639 point predictions, 189 incorrect
(29.58\% conditional error). A history filter chosen from those data
and tested on a separate 1,600-draw sample issued 436 predictions with
94 errors (21.56\%). Both failed the joint target of at most 5\%
conditional error and at least 25\% prediction coverage. The present
certificate comparison has a different estimand.

Reanalysis of the second sample with a covering class found 31 initially
singleton targets among 1,600 records. A response-branch scan on those
records found 134 possible gains and zero possible losses for the
selected policy versus the fixed query, but collected no observations.
It did not establish realized superiority. The policy was chosen with
these results known, then frozen before the new class-frame sample.
Development records are excluded from the new denominator.

\subsection{Class, target and sampling contract}

The candidate class contains deterministic systems with at most three
states and binary inputs and outputs, covered by 46,656 labeled
three-state tables. Labeled candidates and sampled behavior classes
are different units. A certificate requires a nonempty compatible set
unanimous about the specified two-output target from the original
cutoff. Coverage and implementation correctness are premises, not
established for arbitrary worlds by observing no mistakes. Actors receive
the legal history, target definition and candidate class, without true
system tables or state. The observation tool returns only a permitted
intermediate response, never the target continuation.

Write the transition/output table as $(\delta(s,a),\lambda(s,a))$,
with $s\in\{0,1,2\}$, $a\in\{0,1\}$ and initial state 0. The table
emits $\lambda(s,a)$ on a transition from $s$ and moves to
$\delta(s,a)$. The eight observed input--output pairs determine each
compatible candidate's cutoff state $s_8$. Query $q_b$ applies just
input $b$ to a fresh copy of that cutoff and returns
$\lambda(s_8,b)$, for $b\in\{0,1\}$. The target is
\[
T(s_8)=\bigl(\lambda(\delta(s_8,0),0),\,
                 \lambda(\delta(s_8,1),0)\bigr).
\]
These are the second outputs of $(0,0)$ and $(1,0)$, respectively,
starting from separate copies of the original cutoff. The first outputs
are not the target; the chosen query does not shift its reference state.
The stored target code is $2T_0+T_1\in\{0,1,2,3\}$. The fixed arm
always chooses $q_0$; $q_0$ also wins the selected policy's final tie.
Only $q_0$ and $q_1$ are legal: a query cannot request the second step
or the target pair. Actor-visible definitions and candidate-side
computation are disclosed in \texttt{artifact/inspection/wire.py} and
\texttt{compute.py}; they do not supply the true system's table.

Canonicalization gives 11,500 behaviors in the full at-most-three-state
class: 4 with one minimal reachable state, 144 with two, and 11,352 with
three. The original sampling design restricts worlds to exactly three
minimal reachable states, removing the 148 lower-state behaviors. Seven
earlier development fixtures represent two distinct eligible three-state
behaviors; removing these gives the frozen 11,350-class catalog. The first
1,600-draw development sample uses 1,499 distinct classes, leaving 9,851;
the second uses 1,479 additional classes, leaving 8,372. Each has 256
eight-input histories. Thus the apparent 150-class gap is an explicit
sampling restriction and prior-fixture exclusion, not a change to the
actor's covering hypothesis class. The two classes and all subsequent
exclusions are reconstructible by canonical-table identity and index in
\texttt{artifact/acquisition/replay.py}. The original local structural
intent is dated 26 October 2026, before the 27 October main freeze;
these local dates and hashes are not independent public time attestations.
The fixed sample contains 5,000 draws, uniform with replacement over
their 2,143,232 pairs. The implementation uses an operating-system seed,
a SHA-256 counter stream and rejection mapping to avoid index-modulo
bias. Six upper-residue rejections among 5,006 stream words are sampling
implementation events, not outcome exclusions. The inferential model
is independent uniform draws from this finite frame, not a claim that
cryptographic expansion provides physically independent randomness.
All seven duplicate-pair excess draws and 97 initially singleton targets
remain in $N$; 4,903 draws are initially ambiguous. Visible-history
overlap is allowed. Class exclusion does not imply absence from catalog
construction or all prior research.

\subsection{Pre-specified interval, size and decision}

Let $G$ and $L$ count selected-only and fixed-only certificates, and
$D=G+L$. Write $\theta$ for discordance probability and $\beta$ for the
selected-only probability conditional on discordance. The difference
in marginal certificate probabilities is
\[
\Delta=\theta(2\beta-1),\qquad
\widehat\Delta=(G-L)/N.
\]
The primary procedure forms central 97.5\% Clopper--Pearson intervals
for $\theta$ from $D/N$ and for $\beta$ from $G/D$ (full $[0,1]$ when
$D=0$), then takes the minimum and maximum of
$\theta(2\beta-1)$ at their four rectangle corners. Conditional
coverage for the second interval, unconditional coverage for the first
and the union bound give joint coverage of at least 95\% under the
paired-draw model. This is a conservative projection, not the same
procedure as a marginal Wald interval. The numerical implementation
and boundary cases were checked before sampling.

The pre-specified $N$ grid was 1,000 through 6,000 in steps of 1,000.
At planning difference 0.02 and discordance 0.15 (selected-only 0.085,
fixed-only 0.065), $N=5{,}000$ was the first to reach the 90\% power
target: 91.65415\%, versus 83.92156\% at 4,000. Planning used the same
zero-exclusion decision as the primary interval. For positive
discordance, it requires the central 97.5\% conditional interval for
$\beta$ to exclude one half, rather than substituting a 5\% two-sided
test. Power depends on the assumed discordance; 0.15 was not a proven
upper bound. Two percentage points was a research-interest threshold,
not a utility valuation or a claim that the true effect was that large.

One final readout followed all commitments and the target audit.
With $G=90$, $L=1$, the component intervals are
$\theta\in[0.0142214,0.0229103]$ and
$\beta\in[0.9319176,0.9998618]$. The projected interval is
$[0.01228495885,0.02290400123]$. Its lower endpoint exceeds zero,
so the positive-difference claim is supported; its inclusion of 0.02
leaves the interest threshold unresolved. Technical missingness would
retain the fixed denominator rather than trigger replacement draws.
Any incorrect certificate, identity violation or target leak is a
validity failure independent of the confidence interval. There was no
outcome-based sample extension or policy revision.

\subsection{Execution controls, diagnostics and work}

Separate actor processes and state copies implement the two arms.
The controls exercise invalid-query rejection, arm-order invariance
of normalized scientific outputs, masked-target invariance, rejection
of target-continuation access, pre-audit commitments, incorrect-certificate
detection and complete work accounting. Five thousand paired predictions
were committed before privileged target evaluation. The 10,000 actor
processes and keep-awake process exited. Pre-main checks used 77
synthetic extra observations and 81 actor processes; they are excluded
from the scientific sample. Same-host checks and clean-copy replay are
not independent laboratory validation.

Among initially ambiguous draws, the fixed and selected policies
certified 186/4,903 and 275/4,903 targets. Among the 686 draws with
different queries, the counts were 31/686 and 120/686. These are
descriptive conditional rates. The one fixed-only result means that
counting singleton-target branches does not maximize certification for
every actual response. Likewise, the 4,627 neither-certified pairs
do not prove that every permitted one-observation method would fail.

\begin{table}[!ht]
\centering
\small
\begin{tabular}{lrr}
\toprule
Measured operation & Fixed & Selected\\
\midrule
Environment extra observations & 5,000 & 5,000\\
Candidate--history checks & 233,280,000 & 233,280,000\\
Candidate history steps & 494,337,048 & 494,337,048\\
Query--candidate evaluations & 5,283,246 & 10,566,492\\
Selection-branch checks & 0 & 20,000\\
Target--candidate evaluations & 5,283,246 & 5,283,246\\
Target--candidate steps & 21,132,984 & 21,132,984\\
Certificate-branch lookups & 5,000 & 5,000\\
Reused query values within an arm & 0 & 5,283,246\\
Actor CPU seconds & 136.38 & 137.96\\
Actor wall seconds & 143.91 & 146.28\\
Actor-process wall seconds & 311.74 & 313.80\\
\bottomrule
\end{tabular}
\caption{Work for the separate 5,000-draw acquisition study.
Counters are operation-specific and may overlap; reused values are
not additional evaluations. No cross-arm or cross-record caching is
used. Timing does not identify an order-independent speed difference.}
\label{tab:observation-work}
\end{table}


The total main-run elapsed time was 640.28 seconds. In addition to
10,000 extra observations, environment accounting records 40,000
history-generation steps, 5,000 resets, 10,000 clones, zero history-replay
steps and 20,000 target-audit steps. The audit compares 655 issued
certificates to their targets. These operation-specific counters are
not an additive cost score. Timing is descriptive because the fixed arm
always ran first. Model/provider requests were zero; research-assistant
use was unmeasured.


\subsection{Raw-input reader reconstruction}
The bundle includes the original catalogs, exclusions, seed, rejection
receipts, frozen plan, commitments, events, histories and arm records,
plus collector/environment source as inspection-only text. Plain JSON
columns and string references retain every original field and reproduce
all 5,000 original record hashes. Retained process IDs and timings establish byte identity only.

Run \texttt{python3.12 -I -B -S artifact/acquisition/replay.py} from the
bundle root. It reconstructs the universe, ordered catalogs, exclusions
and full first development manifest; verifies the original hashes and
5,006 sampling receipts; regenerates histories, both policies' query
responses, candidate computations and certificates; and fixes predictions
before recomputing targets. Every scientific field and work counter is
compared before recomputing the primary interval. The archived
commitment-before-truth event order is also checked.

This standard-library replay collects no new sample, calls no model, and
executes neither participant programs nor the historical isolation collector.
Independent historical attestation remains outside its scope.

\end{document}
